\chapter*{Introduction générale}
\markboth{INTRODUCTION GÉNÉRALE}{}
\addcontentsline{toc}{chapter}{Introduction générale}

Dans le secteur des biens de grande consommation, la performance commerciale ne dépend pas uniquement de la disponibilité des produits dans les entrepôts ou dans les circuits de distribution. Elle dépend aussi de la qualité de l’exécution réalisée en magasin. La présence effective des produits en rayon, le respect des standards de merchandising, la visibilité des marques, la conformité du placement et la fiabilité des informations remontées depuis le terrain constituent des éléments essentiels pour piloter l’activité commerciale et accompagner la prise de décision.

Dans ce contexte, les données issues des opérations terrain occupent une place importante. Elles permettent de suivre l’activité des merchandisers, d’évaluer la couverture des points de vente, de contrôler la disponibilité des produits, d’identifier les écarts d’exécution et de produire des indicateurs destinés aux équipes internes ainsi qu’au client. Cependant, lorsque ces données sont exploitées quotidiennement à travers plusieurs fichiers Excel, rapports séparés et traitements manuels, leur utilisation devient plus lourde, moins centralisée et plus difficile à suivre dans le temps.

Le présent projet de fin d’études s’inscrit dans le cadre d’un stage réalisé au sein de Smollan Morocco, en tant que Data Analyst Intern, autour du projet Unilever. Dans ce projet, les merchandisers effectuent leurs visites en magasin à travers l’application mobile Smart 3. Les données collectées sont ensuite synchronisées vers le portail Smollan, à partir duquel les équipes opérationnelles exportent principalement deux fichiers : \textit{Data Dump} et \textit{Coverage}. Ces fichiers constituent la base du suivi quotidien, puisqu’ils alimentent les rapports de couverture, de déviation, de disponibilité produit, de qualité d’exécution, de planogramme, de placement et d’autres indicateurs liés au merchandising terrain.

L’exploitation de ces fichiers répond à un besoin opérationnel réel, mais elle présente aussi plusieurs limites. Les données sont dispersées entre plusieurs exports, les traitements sont répétés, l’historique reste difficile à consolider et certaines vérifications nécessitent encore des analyses manuelles. Cette situation peut ralentir la préparation des rapports et compliquer l’identification rapide des situations nécessitant une attention particulière, notamment lorsqu’il s’agit de croiser les informations de couverture, les réponses terrain et les indicateurs métier.

Le sujet de ce PFE porte donc sur la conception d’une solution data-driven pour le suivi et l’analyse des opérations de merchandising terrain du projet Unilever chez Smollan Morocco. L’objectif n’est pas de remplacer totalement les outils existants, mais de proposer une approche complémentaire permettant de mieux structurer les données, d’automatiser une partie des traitements, de centraliser l’historique, d’appliquer des règles de validation métier et de restituer les résultats à travers des supports adaptés aux différents usages.

La solution proposée s’appuie sur une chaîne de traitement des données permettant d’exploiter les fichiers \textit{Data Dump} et \textit{Coverage}, de les transformer et de les charger dans une base MySQL centralisée. Cette base constitue le socle de la solution. Elle alimente ensuite plusieurs composantes : un moteur de validation destiné à produire des signaux de contrôle internes, des services backend permettant d’exposer les données selon les usages, des tableaux de bord Power BI pour l’analyse décisionnelle, une application mobile de consultation côté client et un portail back-office réservé au suivi interne des situations détectées.

La démarche adoptée repose d’abord sur la compréhension du processus métier existant, depuis la visite effectuée par le merchandiser jusqu’à l’exploitation des données par les équipes internes. Elle se poursuit par l’analyse des sources de données, la conception d’une architecture adaptée, le développement du pipeline Data Engineering, la structuration de la base relationnelle, la mise en place des règles de validation et le développement des couches de restitution. Cette progression permet de relier les choix techniques aux besoins opérationnels observés durant le stage.

Le rapport est structuré selon la progression du projet. Il commence par la présentation du contexte général, de l’organisme d’accueil et de la problématique traitée. Il analyse ensuite l’existant et les besoins identifiés, avant de présenter la conception globale de la solution retenue. Les chapitres suivants décrivent la mise en place des principales composantes techniques et applicatives, puis évaluent les résultats obtenus, les apports du projet et les limites de la solution réalisée.